\section{Respuesta a las preguntas de investigación}
\label{sec:disc_preguntas}
% Extensión estimada: 1.5 páginas, cuatro párrafos (ficha 6.1).
% Fuentes: sec:res_clasif_discusion, sec:res_gen_discusion; tab:res_pruebas_vlm, tab:res_dermlip,
% fig:res_melanoma_semillas, tab:res_estado_arte (sensibilidad de melanoma de Swin-Tiny con cada regla),
% tab:res_gen_cascada, tab:res_gen_ablaciones, tab:res_gen_pipeline. Incluye el párrafo 1 del antiguo 4.6.

La primera pregunta de investigación (Sección~\ref{sec:intro_objetivos}) es si los \glspl{vlm} clasifican las
lesiones de piel, y en particular el melanoma, tan bien como los clasificadores entrenados solo con imágenes. En la
prueba de \emph{HAM10000}, ningún \gls{vlm} usado sin una cabeza entrenada sobre ese conjunto lo logra. Los
\glspl{vlm} generativos de propósito general quedan en el nivel del azar con las cinco estrategias evaluadas, y el
\gls{vlm} contrastivo especializado, DermLIP, queda entre 0.15 y 0.18 de sensibilidad macro por debajo de los tres
mejores supervisados, con intervalos de \emph{bootstrap} que excluyen el cero (Sección~\ref{sec:res_comparacion}). En
melanoma la distancia es parecida: DermLIP sin ajuste detecta la mitad de los melanomas, mientras que Swin-Tiny
detecta unos tres de cada cinco con la clase más probable y cuatro de cada cinco con el umbral fijado en validación.

La respuesta tiene un matiz. Con una cabeza entrenada sobre su \emph{embedding}, DermLIP se acerca a los supervisados
en sensibilidad macro, aunque no en F1 macro (Sección~\ref{sec:res_vlm_dermlip}). En ese uso, sin embargo, el
\gls{vlm} aporta la representación y el diagnóstico lo decide un clasificador entrenado con las etiquetas de
\emph{HAM10000}, de modo que la comparación ya no es entre un \gls{vlm} y un clasificador supervisado, sino entre dos
clasificadores supervisados con distinta representación. La respuesta a la primera pregunta es, entonces, negativa
para los \glspl{vlm} usados como tales, y positiva solo para la representación de un \gls{vlm} especializado.

La segunda pregunta es si los reportes que genera un \gls{vlm} afirman el diagnóstico correcto y describen la lesión
que muestra la imagen. En cuanto al diagnóstico, un generador sin anclaje no afirma de manera confiable la clase de la
imagen (Sección~\ref{sec:res_gen_directa}), pero en la cascada el reporte afirma el diagnóstico del clasificador, y la
consistencia diagnóstica de Mini-VLM, 0.849, coincide con la exactitud de Swin-Tiny. Con el ruteo por concordancia,
esa consistencia llega a 0.870 sobre la prueba completa (Sección~\ref{sec:res_gen_pipeline}). En cuanto a la
descripción, la imagen cambia lo que el reporte describe de cada lesión, lo que verifica en parte su fidelidad visual,
pero la descripción también se adapta al diagnóstico recibido (Sección~\ref{sec:res_gen_cualitativos}), y que cada
rasgo descrito esté en la imagen no se verificó. La respuesta a la segunda pregunta es, por lo tanto, afirmativa para
el diagnóstico en la medida en que lo es el clasificador que lo ancla, y solo parcial para la descripción.

Las dos respuestas están relacionadas. Como un \gls{vlm} generativo no distingue las siete clases, el diagnóstico del
reporte no puede salir del generador, y por eso el \emph{pipeline} lo toma de un clasificador supervisado. A su vez,
como el generador transcribe el diagnóstico que recibe, la fidelidad diagnóstica del reporte queda fijada por ese
clasificador. Las secciones siguientes discuten qué explica la primera respuesta, qué implica la segunda para el
diseño de un sistema de generación de reportes y hasta dónde llega lo que miden las métricas.
